\documentclass[t,handout]{beamer}
\usepackage[utf8]{inputenc}
\usepackage[ngermanb]{babel}
\usepackage{listings}
\lstset{ %
  basicstyle=\ttfamily\color{black}\tiny,
  columns=fullflexible,
  breaklines=true,
  frame=single
}

\usepackage{hyperref}
\definecolor{darkblue}{rgb}{0,0,.5}
\hypersetup{pdftex=true, colorlinks=true, breaklinks=true, linkcolor=darkblue, menucolor=darkblue, pagecolor=darkblue, urlcolor=darkblue}

\title{Taskwarrior -- Setup Taskserver}
\subtitle{Your guide to success}
\author[Deimeke, Dirk]{Dirk Deimeke}
\institute[Taskwarrior Academy]{Taskwarrior Academy}
\date{March 2016}
\titlegraphic{\includegraphics[width=3cm,height=3cm]{tw-xxl}}
\subject{Taskwarrior}
\keywords{taskwarrior, task, management, commandline, getting things done}

\setbeamercovered{transparent}

\pgfdeclareimage[width=5mm]{tw-logo}{tw-xl}

%%%%%%%%%%%%%%%%%%%%%%%%%%%
% Layout - Start
%%%%%%%%%%%%%%%%%%%%%%%%%%%

\setbeamertemplate{navigation symbols}{}

\useinnertheme{default}
\useoutertheme{infolines}
\usefonttheme{structurebold}

\newlength{\boxwidth}
\setlength{\boxwidth}{121px}
\setbeamertemplate{headline}{%
\begin{beamercolorbox}[dp=3px,ht=6px,wd=\boxwidth,center]{palette tertiary}%
\insertauthor\ (\insertinstitute)%
\end{beamercolorbox}%
\vskip-9px\hskip\boxwidth
\begin{beamercolorbox}[dp=3px,ht=6px,wd=\boxwidth,center]{palette secondary}%
\inserttitle%
\end{beamercolorbox}%
\begin{beamercolorbox}[dp=3px,ht=6px,wd=\boxwidth,right]{palette primary}%
\insertdate\hskip15px\insertframenumber\,/\,\inserttotalframenumber\hspace*{8px}
\end{beamercolorbox}%
}
\setbeamertemplate{footline}{}

%%%%%%%%%%%%%%%%%%%%%%%%%%%
% Colors - Start
%%%%%%%%%%%%%%%%%%%%%%%%%%%

\definecolor{basecolor}{gray}{0.2}
\definecolor{mittelgrau}{gray}{0.4}
\definecolor{hellgrau}{gray}{0.93}
\definecolor{gelb}{rgb}{1.0,1.0,0.75}

% infolines color
\setbeamercolor{palette primary}{fg=white,bg=mittelgrau}
\setbeamercolor{palette secondary}{fg=white,bg=basecolor}
\setbeamercolor{palette tertiary}{fg=white,bg=black}

% frame and title color
\setbeamercolor{frametitle}{fg=white,bg=basecolor}
\setbeamercolor{titlelike}{fg=white,bg=basecolor}
\setbeamerfont{frametitle}{series=\bfseries}

% TOC color
\setbeamercolor{section in toc}{fg=basecolor}

% text color
\setbeamercolor{normal text}{fg=black,bg=white}

% item color
\setbeamercolor{item}{fg=basecolor}
\setbeamercolor{itemize item}{fg=black}

% block color
\setbeamercolor{block title}{fg=black,bg=gelb}
\setbeamercolor{block body}{fg=black,bg=gelb}

% block alert color
\setbeamercolor{block title alerted}{fg=black,bg=gelb}
\setbeamercolor{block body alerted}{fg=black,bg=gelb}

% block example color
\setbeamercolor{block title example}{fg=black,bg=gelb}
\setbeamercolor{block body example}{fg=black,bg=gelb}

\setbeamertemplate{frametitle}{%
\vskip-1px%
\begin{beamercolorbox}[wd=363px,ht=20px,dp=12px]{frametitle}
\usebeamerfont{frametitle}%
\usebeamercolor{frametitle}%
\vspace*{-5px}\hspace*{10px}\includegraphics[width=24px,height=24px]{tw-xl.png}\\
\vspace*{-20px}\hspace*{40px}\insertframetitle%
\end{beamercolorbox}
}

%%%%%%%%%%%%%%%%%%%%%%%%%%%
% Colors - Stop
%%%%%%%%%%%%%%%%%%%%%%%%%%%

\begin{document}

%\begin{frame} % Titel
%	\titlepage
%\end{frame}

% \logo{\pgfuseimage{tw-logo}}

%\begin{frame}\frametitle{Content}
%	\tableofcontents
%\end{frame}

%%%%%%%%%%%%%%%%%%%%%%%%%%%
\section{Preparation}
%%%%%%%%%%%%%%%%%%%%%%%%%%%

\parskip1ex

\begin{frame}[fragile]\frametitle{Preparation -- Backup Your Data}\label{backup}

    Let's reinforce a good habit and make a backup copy of your data first. Here is a very easy way to backup your data:

    \begin{lstlisting}
        $ cd ~/.task
        $ tar czf task-backup-$(date +'%Y%m%d').tar.gz *
    \end{lstlisting}

    Now move that file somewhere safe. All software contains bugs, so make regular backups.

	\begin{alertblock}{Attention!}
		This is not only due to a good habit, we will modify your data, so this backup is highly recommended.
	\end{alertblock}

\end{frame}

\begin{frame}[fragile]\frametitle{Preparation -- Choose A  Machine}

	A suitable machine to run your Taskserver is one that is always available. If you have such a machine, or have access to a hosted machine, that is ideal.

	If your machine is not continuously available, it can still be a suitable Taskserver because the sync mechanism doesn't require continuous access. When a client cannot sync, it simply accumulates local, unpropagated changes until it can sync.

	A laptop is a poor choice for a Taskserver host.

\end{frame}


\begin{frame}[fragile]\frametitle{Preparation -- Choose A  Port}

	By default, Taskserver uses port 53589. You can choose any port you wish, provided it is unused. If you choose a port number that is under 1024, then Taskserver must run as root, which is not recommended.

\end{frame}

\begin{frame}[fragile]\frametitle{Preparation -- User/Group}
  Ideally you will create a new user and group solely to run the
  Taskserver. This helps you keep the data secure from other users
  on the machine, as well as controlling the privileges of
  Taskserver.
\end{frame}

\begin{frame}[fragile]\frametitle{Preparation -- Firewall}
  Depending on what devices you use to access your server, you may
  need to configure the firewall to allow incoming TCP/IP traffic on
  your chosen port.
\end{frame}


%%%%%%%%%%%%%%%%%%%%%%%%%%%
\section{Installation}
%%%%%%%%%%%%%%%%%%%%%%%%%%%

\iffalse
<a name="installation"></a>
<h3><span class="label label-success">2.1</span> Installation</h3>
<p>
  Installing Taskserver from a binary package is the simplest
  option, but you will need to refer to your package manager's
  documentation and procedures for doing this.
</p>

<p>
  Take a look at the
  <a href="/download/#dist">Download</a> page
  for examples.  Generally there are too many package managers to
  make a complete list with instructions here.
</p>

<p>
  Most importantly, for now, Taskserver is a new product, and there
  are very few packages available.  It is expected that this
  situation will change soon.  When it does, this page will be updated.
</p>

<p>
  Meanwhile, there is installation from either git or tarball.
</p>

%%%

<h3><span class="label label-success">2.2</span> Installation</h3>
<p>
  Installing Taskserver from a tarball is a matter of downloading
  the tarball, extracting it, satisfying dependencies and building
  the server.
</p>

<a name="depends"></a>
<h4>Dependencies</h4>
<p>
  Before building the software, you will need to satisfy the
  dependencies by installing the following:

  <ul>
    <li>GnuTLS (likely <code>libgnutls-dev</code>, ideally version 3.2 or newer)</li>
    <li><code>libuuid</code></li>
    <li>CMake (2.8 or newer)</li>
    <li><code>make</code></li>
    <li>A C++ Compiler (GCC 4.7 or Clang 3.0 or newer)</li>
  </ul>

  Note that some OSes (Darwin, FreeBSD ...) include <code>libuuid</code>
  functionality in libc.
</p>

<p>
  You don't necessarily need the latest version of all components,
  but it is a good idea if you can.  GnuTLS is a security component,
  and as such, it is very important that it is current.

  Using GnuTLS version 2.12.x is neither adequately secure, nor
  production quality.
</p>

<a name="download"></a>
<h4>Download</h4>
<p>
  The next step is to obtain the code. This means getting the
  Task Server 1.0.0 (or newer) source tarball.  You should check
  for the latest stable release here:
</p>

<p>
  <a href="http://taskwarrior.org/download/">http://taskwarrior.org/download/</a>
</p>

<p>
  You can download the tarball with <code>curl</code>, as an example
  of just one of many ways to download the tarball.
</p>

<pre>$ curl -O http://taskwarrior.org/download/taskd-latest.tar.gz</pre>

<a name="build"></a>
<h4>Build</h4>
<p>
  Expand the tarball, and build the Taskserver.
</p>

<pre>$ tar xzf taskd-latest.tar.gz
$ cd taskd-latest
$ cmake -DCMAKE_BUILD_TYPE=release .
...
$ make
...</pre>

<a name="install"></a>
<h4>Install</h4>
<p>
  Now install Taskserver.  This copies files into the right place,
  and installs man pages.
</p>

<pre>$ sudo make install
...</pre>

<a name="verify"></a>
<h4>Verify</h4>
<p>
  Run the <code>taskd</code> command to verify that the server is
  installed, and the location is in your <code>$PATH</code>.
  You should see something like this:
</p>

<pre>$ taskd

Usage: taskd -v|--version
taskd -h|--help
taskd diagnostics
taskd validate &lt;JSON | file&gt;
taskd help [&lt;command&gt;]

Commands run only on server:
taskd add     [options] org   &lt;org&gt;
taskd add     [options] group &lt;org&gt; &lt;group&gt;
taskd add     [options] user  &lt;org&gt; &lt;user&gt;
taskd config  [options] [--force] [&lt;name&gt; [&lt;value&gt;]]
taskd init    [options]
taskd remove  [options] org   &lt;org&gt;
taskd remove  [options] group &lt;org&gt; &lt;group&gt;
taskd remove  [options] user  &lt;org&gt; &lt;user&gt;
taskd resume  [options] org   &lt;org&gt;
taskd resume  [options] group &lt;org&gt; &lt;group&gt;
taskd resume  [options] user  &lt;org&gt; &lt;user&gt;
taskd server  [options] [--daemon]
taskd status  [options]
taskd suspend [options] org   &lt;org&gt;
taskd suspend [options] group &lt;org&gt; &lt;group&gt;
taskd suspend [options] user  &lt;org&gt; &lt;user&gt;

Commands run remotely:
taskd client  [options] &lt;host:port&gt; &lt;file&gt; [&lt;file&gt; ...]

Common Options:
--quiet        Turns off verbose output
--debug        Generates debugging diagnostics
--data &lt;root&gt;  Data directory, otherwise $TASKDDATA
--NAME=VALUE   Temporary configuration override
</pre>

%%%

<h3><span class="label label-success">2.3</span> Installation</h3>
<p>
  Installing Taskserver from git is a matter of cloning the git
  repository and building the server.
</p>

<a name="depends"></a>
<h4>Dependencies</h4>
<p>
  Before building the software, you will need to satisfy the
  dependencies by installing the following:

  <ul>
    <li><code>git</code></li>
    <li>GnuTLS (likely <code>libgnutls-dev</code>, ideally version 3.2 or newer)</li>
    <li><code>libuuid</code></li>
    <li>CMake (2.8 or newer)</li>
    <li><code>make</code></li>
    <li>A C++ Compiler (GCC 4.7 or Clang 3.0 or newer)</li>
  </ul>

  Note that some OSes (Darwin, FreeBSD ...) include <code>libuuid</code>
  functionality in libc.
</p>

<p>
  You don't necessarily need the latest version of all components,
  but it is a good idea if you can.  GnuTLS is a security component,
  and as such, it is very important that it is current.

  Using GnuTLS version 2.12.x is neither adequately secure, nor
  production quality.
</p>

<a name="clone"></a>
<h4>Clone</h4>
<p>
  Now clone the repository like this:
</p>

<pre>$ git clone https://git.tasktools.org/scm/tm/taskd.git taskd.git
...</pre>

<div class="alert alert-info">
  <p>
    It is hightly recommended that you build the stable version.
    This involves simply moving on to the next step,
    <a href="#build">'Build'</a>.
    Only under special circumstances should you build the unstable
    development version.
  </p>

  <p>
    The unstable development version is at no point guaranteed to
    work or even compile. The only time it does ѕtabilize is right at
    the end of the development cycle, and in that case, you should
    wait until the release, so you are running a supported version.
  </p>

  <p>
    The stable version is always merged to the <code>master</code>
    branch, which is the default branch, so ordinarily nothing needs
    to be done. To build an unstable branch, first determine which
    branch by looking at the available branches:
  </p>

  <pre>$ cd taskd.git
$ git branch -a
* master
remotes/origin/1.1.0
remotes/origin/1.1.1
remotes/origin/1.2.0
remotes/origin/HEAD -&gt; origin/master
remotes/origin/master
</pre>

  <p>
    The convention we use is that <code>master</code> represents the
    stable release.  The numbered branches represent the latest
    development (1.2.0, the 'highest' branch number, ending in '.0')
    and a patch branch (1.1.1, ending in a non-zero number).
  </p>

  <p>
    Patch branches are reserved for emeergency releases, so in this
    example you would choose to build 1.2.0 as the latest development
    branch like this:
  </p>

  <pre>$ git checkout 1.2.0
Branch 1.2.0 set up to track remote branch 1.2.0 from origin.
Switched to a new branch '1.2.0'</pre>
</div>

<a name="build"></a>
<h4>Build</h4>
<p>
  Now build the Taskserver.
</p>

<pre>$ cd taskd.git
$ cmake -DCMAKE_BUILD_TYPE=release .
...
$ make
...</pre>

<a name="test"></a>
<h4>Test</h4>
<p>
  Having built the server, now build and run the unit tests.
  Although this is an optional step, it is a good idea to know
  whether the build works on your platform.
</p>

<pre>$ cd test
$ make
...
$ ./run_all

Pass:     2920
Fail:        0
Skipped:     0
Runtime:     1 seconds

$ cd ..
$</pre>

<p>
  This example shows that all 2,920 tests pass.  If you see test
  failures, stop and report them.

  Note that there are some unit tests that fail if you have not
  built the latest commit. Seeing 4 test failures may mean all is
  well.  Seeing 30 failures does not.
</p>

<a name="install"></a>
<h4>Install</h4>
<p>
  Now install Taskserver.  This copies files into the right place,
  and installs man pages.
</p>

<pre>$ sudo make install
...</pre>

<a name="verify"></a>
<h4>Verify</h4>
<p>
  Run the <code>taskd</code> command to verify that the server is
  installed, and the location is in your <code>$PATH</code>.
  You should see something like this:
</p>

<pre>$ taskd

Usage: taskd -v|--version
taskd -h|--help
taskd diagnostics
taskd validate &lt;JSON | file&gt;
taskd help [&lt;command&gt;]

Commands run only on server:
taskd add     [options] org   &lt;org&gt;
taskd add     [options] group &lt;org&gt; &lt;group&gt;
taskd add     [options] user  &lt;org&gt; &lt;user&gt;
taskd config  [options] [--force] [&lt;name&gt; [&lt;value&gt;]]
taskd init    [options]
taskd remove  [options] org   &lt;org&gt;
taskd remove  [options] group &lt;org&gt; &lt;group&gt;
taskd remove  [options] user  &lt;org&gt; &lt;user&gt;
taskd resume  [options] org   &lt;org&gt;
taskd resume  [options] group &lt;org&gt; &lt;group&gt;
taskd resume  [options] user  &lt;org&gt; &lt;user&gt;
taskd server  [options] [--daemon]
taskd status  [options]
taskd suspend [options] org   &lt;org&gt;
taskd suspend [options] group &lt;org&gt; &lt;group&gt;
taskd suspend [options] user  &lt;org&gt; &lt;user&gt;

Commands run remotely:
taskd client  [options] &lt;host:port&gt; &lt;file&gt; [&lt;file&gt; ...]

Common Options:
--quiet        Turns off verbose output
--debug        Generates debugging diagnostics
--data &lt;root&gt;  Data directory, otherwise $TASKDDATA
--NAME=VALUE   Temporary configuration override
</pre>


%%%%%%%%%%%%%%%%%%%%%%%%%%%
\section{Server Setup}
%%%%%%%%%%%%%%%%%%%%%%%%%%%

<h3><span class="label label-success">3</span> Server Configuration</h3>
<p>
  Configuring the server is straightforward, but needs a little
  planning.
</p>

<a name="location"></a>
<h4>Choose a Data Location</h4>
<p>
  A location for the data must be chosen and created.  The
  <code>TASKDDATA</code> environment variable will be used to
  indicate that location to all the <code>taskd</code> commands.
</p>

<pre>$ export TASKDDATA=/var/taskd
$ mkdir -p $TASKDDATA</pre>

<p>
  If the <code>TASKDDATA</code> variable is not set, then most
  <code>taskd</code> commands require the <code>--data ...</code>
  argument, otherwise the commands rely on the <code>TASKDDATA</code>
  value to indicate the location.

  Everything the server does will be confined to that directory.
</p>

<p>
  Note that there are two 'D's in <code>TASKDDATA</code>, and
  omitting one is a common mistake to make.
</p>

<p>
  Note that there must be write permission in the directory for the
  user that will run the server.
</p>

<a name="init"></a>
<h4>Initialization</h4>
<p>
  Now we let the server initialize that directory:
</p>

<pre>$ taskd init
You must specify the 'server' variable, for example:
taskd config server localhost:53589

Created /var/taskd</pre>

<a name="keys"></a>
<h4>Keys & Certificates</h4>
<p>
  Now we create certificates and keys. The command below will
  generate all the certs and keys for the server, but this uses
  self-signed certificates, and this is not recommended for
  production use. This is for personal use, and this may be
  acceptable for you, but if not, you will need to purchase a
  proper certificate and key, backed by a certificate authority.
</p>

<p>
  The certificate and key generation scripts make assumptions
  <b><em>that are guaranteed to be wrong for you</em></b>.
  Specifically the
  <code>generate.server</code> script has a hard-coded
  <code>CN</code> entry that is not going to work.  You
  <b><em>need</em></b> to edit the <code>vars</code> file.
</p>

<pre>...
CN=localhost
...</pre>

<p>
  You will need to modify this value to match your server.
  It is this value against which Taskwarrior validates the server
  name, so use a value similar to <code>ack.tasktools.org</code>,
  which is what we use, but of course don't expect that to work for
  you. If you do not change this value, the only option for the
  client is to skip some or all certificate validation,
  <b><em>which is a bad idea.</em></b>
</p>

<p>
  Go to your <code>taskd</code> directory, which depends on
  which installation method you chose. Here is is assumed that
  you installed from the source tarball.
</p>

<pre>$ cd ~/taskd-1.1.0/pki
$ ./generate
...
$
$ cp client.cert.pem $TASKDDATA
$ cp client.key.pem  $TASKDDATA
$ cp server.cert.pem $TASKDDATA
$ cp server.key.pem  $TASKDDATA
$ cp server.crl.pem  $TASKDDATA
$ cp ca.cert.pem     $TASKDDATA
$
$ taskd config --force client.cert $TASKDDATA/client.cert.pem
$ taskd config --force client.key $TASKDDATA/client.key.pem
$ taskd config --force server.cert $TASKDDATA/server.cert.pem
$ taskd config --force server.key $TASKDDATA/server.key.pem
$ taskd config --force server.crl $TASKDDATA/server.crl.pem
$ taskd config --force ca.cert $TASKDDATA/ca.cert.pem</pre>

<p>
  There are three classes of key/cert here. There is the CA
  (Certificate Authority) cert, which has cert signing capabilities
  and is used to sign and verify the other certs.

  There are the server key/certs, which are used to authenticate the
  server and encrypt.

  Finally there are client key/certs, which are not what you might
  expect. These are for API access, and not for your Taskwarrior
  client. Those are created later.
</p>

<a name="other"></a>
<h4>Configuration</h4>
<p>
  Now we configure some basic details for the server.  The chosen
  port is 53589. Note that we allow Taskwarrior clients
  specifically.
</p>

<pre>$ cd $TASKDDATA/..
$ taskd config --force log $PWD/taskd.log
$ taskd config --force pid.file $PWD/taskd.pid
$ taskd config --force server localhost:53589</pre>

<p>
  Note that we have chosen <code>localhost:53589</code>, but this
  choice has consequences. The name <code>localhost</code> is not
  network visible, which limits the server to only serving clients
  on the same machine. Use your full machine name for proper network
  addressability.
</p>

<p>
  You can look at all the configuration settings:
</p>

<pre>$ taskd config</pre>

<p>
  You can view all the supported settings with:
</p>

<pre>$ man taskdrc</pre>
</div>

<br />
<br />

%%%

<a name="control"></a>
<h3><span class="label label-success">3</span> Server Start/Stop</h3>
<p>
  You can now to launch the server:
</p>

<pre>$ taskdctl start</pre>

<p>
  This command launched the server as a daemon process. This command
  requires the <code>TASKDDATA</code> variable.  Your server is now
  running, and ready for syncing.

  Note that to stop the server, you use:
</p>

<pre>$ taskdctl stop</pre>

<p>
  Check that your server is running by looking in the <code>taskd.log</code>
  file, or running this:
</p>

<pre>$ ps -leaf | grep taskd</pre>

<a name="interactive"></a>
<h4>Interactive or Non-Daemon Server</h4>
<p>
  A daemon server is typically how you would want to run Taskserver,
  but there may be times when you need to run the server attached to
  a terminal.  These two commands are identical:
</p>

<pre>$ taskdctl start
$ taskd server --data $TASKDDATA --daemon</pre>

<p>
  By omitting the <code>--daemon</code> option, the server remains
  attached to the terminal.  Then to stop the server you can enter
  <code>Ctrl-C</code>.
</p>

<p>
  The interactive mode is really only useful for debugging, in
  conjunction with TLS debug mode, like this:
</p>

<pre>$ taskd config debug.tls 3
$ taskd server --data $TASKDDATA
...</pre>

<p>
  With a <code>debug.tls</code> setting that is non-zero, you see
  lots of security-related diagnostic output.
</p>

<a name="systemd"></a>
<h4>Configure Taskserver to run with a systemd-unit-file</h4>

<p>
You can start Taskserver using a systemd-unitfile like the following (please add the contents of <code>$TASKDDATA</code> not the variable itself). Running the Taskserver as root is not recommended, please add an appropriate user and group to run the daemon with (<code>$TASKDUSER</code> and <code>$TASKDGROUP</code>).
</p>

<pre>[Unit]
Description=Secure server providing multi-user, multi-client access to Taskwarrior data
Requires=network.target
After=network.target
Documentation=http://taskwarrior.org/docs/#taskd

[Service]
ExecStart=/usr/local/bin/taskd server --data $TASKDDATA
Type=simple
User=$TASKDUSER
Group=$TASKDGROUP
WorkingDirectory=$TASKDDATA
PrivateTmp=true
InaccessibleDirectories=/home /root /boot /opt /mnt /media
ReadOnlyDirectories=/etc /usr

[Install]
WantedBy=multi-user.target</pre>

<p>
Afterwards prepare systemd to recognise the file.
</p>

<pre>$ cp taskd.service /etc/systemd/system
$ systemctl daemon-reload
$ systemctl start taskd.service
$ systemctl status taskd.service</pre>

In case everything is running fine, enable the script to start Taskserver on every boot.

<pre>$ systemctl enable taskd.service</pre>

</div>

<br />
<br />


%%%%%%%%%%%%%%%%%%%%%%%%%%%
\section{Client Setup}
%%%%%%%%%%%%%%%%%%%%%%%%%%%

<a name="add"></a>
<h3><span class="label label-success">4</span> Add User to Server</h3>
<p>
  A user account must be created, along with a key, cert and ID,
  before syncing may occur.
</p>

<a name="org"></a>
<h4>Create Organization</h4>
<p>
  Before creating a user account, you may need to create an
  organization. An organization consists of a group of zero or more
  users. You can get away with just one organization, and in this
  example, we will create just one, named 'Public'.
</p>

<p>
  You can create as many organizations as you wish (even one per
  user), and the purpose is simply to group users together.
  Future features will utilize this.
</p>

<pre>$ taskd add org Public
Created organization 'Public'</pre>

<p>
  Now the organization 'Public' exists, we can add users to that
  organization.
</p>

<a name="user"></a>
<h4>Create User</h4>
<p>
  Now we add a new user, named 'First Last' as an example.  You can
  use any name you wish, and if it contains spaces, quote the name
  as shown.
</p>

<pre>$ taskd add user 'Public' 'First Last'
New user key: cf31f287-ee9e-43a8-843e-e8bbd5de4294
Created user 'First Last' for organization 'Public'</pre>

<p>
  Note that you will get a different 'New user key' than is shown
  here, and you will need to retain it, to be used later for client
  configuration.  Note that the key is just a unique id, because
  your name alone is not necessarily unique.
</p>

<a name="cert"></a>
<h4>Create Cert/Key</h4>
<p>
  Go to your <code>taskd</code> directory, which depends on
  which installation method you chose. Here it is assumed that
  you installed from the source tarball.
</p>

<pre>$ cd ~/taskd-1.1.0/pki
$ ./generate.client first_last
...</pre>

<p>
  This will generate a new key and cert, named
  <code>first_last.cert.pem</code>
  and
  <code>first_last.key.pem</code>.
  It is not important that 'first_last' was used here, just that it
  is something unique, and valid for use in a file name. It has no
  bearing on security.
</p>

<a name="client"></a>
<h4>Client</h4>
<p>
  You have now created a new user account on the server, created a
  new client cert and key, and have details that need to be
  transferred to the user, to set up a sync client.  The details
  needed are:

  <ul>
    <li>
      <code>ca.cert.pem</code> is the certificate authority, and
      the only way to validate self-signed certs such as the ones
      we have created here.
    </li>
    <li>
      <code>first_last.cert.pem</code> is the client certificate.
    </li>
    <li>
      <code>first_last.key.pem</code> is the client key.
    </li>
    <li>
      The new user key (yours will be different):
      <code>cf31f287-ee9e-43a8-843e-e8bbd5de4294</code>
    </li>
    <li>
      The organization, <code>Public</code>.
    </li>
    <li>
      The full and proper user name, <code>First Last</code>.
    </li>
    <li>
      The server address and port, <code>host.domain:53589</code>.
      In the
      <a href="/docs/taskserver/configure.html">server configuration</a>
      we used <code>localhost</code> as an example. Whatever you
      actually used there, should be used here.
    </li>
  </ul>
</p>
</div>

<br />
<br />

%%%

<a name="config"></a>
<h3><span class="label label-success">4</span> Configure Taskwarrior</h3>
<p>
  If you have configured Taskserver and created a user account (or
  better yet, someone created an account for you) then you now have
  details needed in the configuration of your Taskwarrior client.
  You should have these files and this information:

  <ul>
    <li>
      CA <code>ca.cert.pem</code>
    </li>
    <li>
      Client certificate: <code>first_last.cert.pem</code>
    </li>
    <li>
      Client key: <code>first_last.key.pem</code>
    </li>
    <li>
      User key (yours will be different): <code>cf31f287-ee9e-43a8-843e-e8bbd5de4294</code>
    </li>
    <li>
      Organization: <code>Public</code>
    </li>
    <li>
      Account name: <code>First Last</code>
    </li>
    <li>
      Server:port <code>host.domain:53589</code>.
      In the
      <a href="/docs/taskserver/configure.html">server configuration</a>
      we used <code>localhost</code> as an example. Whatever you
      actually used there, should be used here.
    </li>
  </ul>

  Now we feed this information to Taskwarrior.
</p>

<p>
  Copy the Cert, Key and CA to your ~/.task directory.  The reason
  we are copying the CA cert is because this is a self-signed cert,
  and we need the CA to validate against.  Alternately we could
  force Taskwarrior to trust all certs, but that is not recommended.
</p>

<pre>$ cp first_last.cert.pem ~/.task
$ cp first_last.key.pem  ~/.task
$ cp ca.cert.pem         ~/.task</pre>

<p>
  Now we need to make Taskwarrior aware of these certs:
</p>

<pre>$ task config taskd.certificate -- ~/.task/first_last.cert.pem
$ task config taskd.key         -- ~/.task/first_last.key.pem
$ task config taskd.ca          -- ~/.task/ca.cert.pem</pre>

<p>
  Now set the server info:
</p>

<pre>$ task config taskd.server      -- host.domain:53589</pre>

<p>
  Finally we provide the credentials, which combine the
  organization, account name and user key:
</p>

<pre>$ task config taskd.credentials -- Public/First Last/cf31f287-ee9e-43a8-843e-e8bbd5de4294</pre>

<div class="callout callout-danger">
  <h4>Trust</h4>
  <p>
    It is possible to configure Taskwarrior's trust level, which
    determines how the server certificate is treated.
  </p>

  <p>
    For Taskwarrior <span class="label label-success">2.3.0</span>
    you can specify <code>taskd.trust=yes</code> in order to skip
    certificate validation.
    <b><em>This is a bad idea.</em></b>
    The default is <code>taskd.trust=no</code>, which does not trust
    the server certificate, which is more secure.
  </p>

  <p>
    For Taskwarrior <span class="label label-success">2.4.0</span>
    you must specify <code>taskd.trust=ignore hostname</code> in
    order to skip certificate hostname validation.
    <b><em>This is a bad idea.</em></b>
    You can also specify <code>taskd.trust=allow all</code> to
    perform no validation.
    <b><em>This is a worse idea.</em></b>
    The default value is <code>taskd.trust=strict</code> which
    performs the most stringent verification, and is more secure.
  </p>
</div>

<p>
  Your Taskwarrior is now ready to sync.
</p>
</div>

<br />
<br />

%%%%%%%%%%%%%%%%%%%%%%%%%%%
\section{Sync}
%%%%%%%%%%%%%%%%%%%%%%%%%%%

<a name="sync"></a>
<h3><span class="label label-success">5</span> Syncing Taskwarrior</h3>
<p>
  You are now ready to sync your Taskwarrior client. You will do this
  differently depending on whether this is the first sync per device,
  or one of the many subsequent syncs.
</p>

<a name="first"></a>
<h4>First Time Sync</h4>
<p>
  The first time you sync is special - the client sends all your
  pending tasks to the server.  This is something you should only
  do once.  Run this:
</p>

<pre>$ task sync init
Please confirm that you wish to upload all your pending tasks to the Task Server (yes/no) yes
Syncing with host.domain:53589

Sync successful.  2 changes uploaded.</pre>

<p>
  You should get an indication that tasks were uploaded, in this
  case 2 of them.
</p>

<a name="general"></a>
<h4>General Sync</h4>
<p>
  After the first time sync, you switch and just use this command:
</p>

<pre>$ task sync
Syncing with host.domain:53589

Sync successful.  No changes.</pre>

<p>
  This will give you feedback about what happened. Please note
  that it is perfectly safe to run this command as often as you
  wish. Syncing is safe and does not consume great system
  resources.
</p>

<p>
  Note that if your client is a mobile device, a sync command may
  consume some of your data usage. Act accordingly.
</p>

<p>
  But it does require network connectivity, and if there is no
  connectivity you will be notified.  It is not a problem if a sync
  fails because of this, because the next sync that works will
  catch up with all the changes, and do the right merging.
</p>

<p>
 Taskwarrior and Taskserver were designed to work together, and
 tolerate intermittent connectivity.
</p>

<a name="reminder"></a>
<h4>Sync Reminder</h4>
<p>
  After you modify data locally, Taskwarrior will start notifying
  you that you need to sync, after commands, like this:
</p>

<pre>$ task project:foo list
No matches.
There are local changes.  Sync required.</pre>

<p>
  This is just a reminder to sync. Respond with a sync, and the
  reminder goes away:
</p>

<pre>$ task sync
Syncing with <server>:<port>

Sync successful.  1 changes uploaded.</pre>

<p>
  If you do not respond with a sync, then local changes accumulate
  unseen by other clients. When you do eventually sync, the data
  will be properly propagated, so it is a question of whether you
  <em>need</em> current data on the server. It is perfectly fine
  to allow <em>weeks</em> to go by without a sync.
</p>
</div>

<br />
<br />

%%% Troubleshooting

<h3>Troubleshooting Sync</h3>
<p>
  Here is a list of problems you may encounter, with the most common
  ones listed first.

  The single most common problem has been that the
  <a href="/docs/taskserver/setup.html">Taskserver Setup Instructions</a>
  were not properly followed.  Please review the steps you took.
</p>

<p>
  It is always a good idea to make sure that you are using the latest
  release of Taskwarrior and Taskserver, not just because bugs are
  fixed that may help you, but also because the solutions below are
  geared toward the current releases.
</p>

<p>
  If you upgrade from an older release of Taskserver, you will need
  to follow the
  <a href="/docs/taskserver/upgrade.html">upgrade instructions</a>.
</p>

<a name="problems"></a>
<h4>Problems</h4>
<table class="table table-striped">
  <tr>
    <td>
      <dt>
        <span class="glyphicon glyphicon-question-sign"></span>
        You tried <code>task sync</code> but Taskwarrior showed you a task list instead
      </dt>
      <dd>
        You have a version of Taskwarrior older than <code>2.3.0</code>,
        which means there was no <code>sync</code> command, and
        you are seeing a list filtered by the search term 'sync'.
        Upgrading is the only solution.
      </dd>
    </td>
  </tr>
  <tr>
    <td>
      <dt>
        <span class="glyphicon glyphicon-question-sign"></span>
        You tried <code>task sync</code> and saw 'Taskwarrior was built without GnuTLS support.  Sync is not available.'
      </dt>
      <dd>
        <p>
          You are using version <code>2.3.0</code> or later, but
          the Taskwarrior binary was compiled without
          <a href="http://www.gnutls.org">GnuTLS</a> support.
        </p>

        <p>
          If you installed Taskwarrior using your OS's package
          manager, you may be suffering from an out of date package.
          Prod your OS's package maintainer for an update.
        </p>

        <p>
          Recent releases make GnuTLS support opt-out instead of
          opt-in, so upgrading to the latest version may help.
          Otherwise, you will need to build Taskwarrior from the
          <a href="/download/task-latest.tar.gz">latest source tarball</a>,
          following the instructions in the <code>INSTALL</code> file.
          If you are a developer, do that. If you are not, then
          installing a development environment is probably not
          something you want to do, in which case contact your
          OS's package maintainer.
        </p>

        <p>
          Verify that your Taskwarrior was built with GnuTLS support
          by running <code>task diagnostics</code>:

          <pre>$ task diagnostics | grep libgnutls
libgnutls: 3.3.18</pre>
        </p>
      </dd>
    </td>
  </tr>

  <tr>
    <td>
      <dt>
        <span class="glyphicon glyphicon-question-sign"></span>
        nodename nor servname provided, or not known
      </dt>
      <dd>
        <p>
          Despite the terrible wording, this means the Taskwarrior
          setting
          <code>taskd.server=&lt;host&gt;:&lt;port&gt;</code>
          refers to a host name that cannot be seen by Taskwarrior.
        </p>
        <p>
          Is it spelled correctly?
          Is the domain correct?
          Is there a valid DNS resolution for the name?
          Is there a firewall between Taskwarrior and Taskserver that
          is not letting through <code>&lt;port&gt;</code> traffic?
        </p>
      </dd>
    </td>
  </tr>

  <tr>
    <td>
      <dt>
        <span class="glyphicon glyphicon-question-sign"></span>
        Could not connect to &lt;host&gt; &lt;port&gt;
      </dt>
      <dd>
        Taskserver may not be running on <code>&lt;host&gt;</code>.
        Check with <code>ps -leaf | grep taskd</code>.
      </đd>
    </td>
  </tr>

  <tr>
    <td>
      <dt>
        <span class="glyphicon glyphicon-question-sign"></span>
        Unable to use port 53589?
      </dt>
      <dd>
        <p>
          By default, port <code>53589</code> is used, but whichever
          you chose must be open on the server.
        </p>

        <p>
          If you are unable to open firewall ports, you can use an SSH
          Tunnel to route port <code>53589</code> traffic over port
          <code>22</code>:

          <pre>$ ssh -L localport:dsthost:dstport user@example.com</pre>
        </p>
      </đd>
    </td>
  </tr>

  <tr>
    <td>
      <dt>
        <span class="glyphicon glyphicon-question-sign"></span>
        Certificate fails validation, Handshake failed
      </dt>
      <dd>
        <p>
          There are many reasons that the TLS handshake can fail.
        </p>

        <p>
          When you generated certificates, you modified a <code>vars</code>
          file, in particular the <code>CN=&lt;name&gt;</code>
          setting. That name must match the output of
          <code>$ hostname -f</code>
          on the server for the certificate to validate.
        </p>

        <p>
          Additionally, that name must also be used in the
          <code>taskd.server=&lt;host&gt;:&lt;port&gt;</code>
          setting for Taskwarrior. Otherwise you'll need
          <code>taskd.trust=ignore hostname</code>.
        </p>

        <p>
          If you are using a self-signed certificate, did you specify
          it using the <code>taskd.ca</code> setting?
        </p>

        <p>
          Setting <code>taskd.ciphers</code> can force the use of
          different ciphers. Use <code>gnutls-cli --list</code> to see
          a list of installed ciphers, and confirm that there is overlap
          between client and server. There needs to be a cipher that
          is available to both, otherwise they cannot communicate.
        </p>
      </đd>
    </td>
  </tr>

  <tr>
    <td>
      <dt>
        <span class="glyphicon glyphicon-question-sign"></span>
        Is your certificate still valid?
      </dt>
      <dd>
        <p>
          Certificates have expiration dates, and if you followed our
          instructions, you created a certificate that is valid for
          one year.  Check your certificate with this command:
          <pre>$ certtool -i --infile ~/.task/&lt;YOUR NAME&gt;.cert.pem</pre>
        </p>

        <p>
          If your certificate has expired, you need a new one.  You
          may also need to regenerate expired server certificates.
        </p>

        <p>
          Note that creating certificates that never expire is a bad
          idea. Certificates may be compromised. A certificate that is
          considered secure today, may not be considered secure in
          a year. Is the key length you chose something that will
          remain suitable in the future? Will the algorithms you
          chose remain secure? For these reasons, choose an
          expiration date that lets you reevaluate your choices in
          the relatively near future.
        </p>
      </đd>
    </td>
  </tr>

  <tr>
    <td>
      <dt>
        <span class="glyphicon glyphicon-question-sign"></span>
        Is your GnuTLS library up to date?
      </dt>
      <dd>
        <p>
          As a
          <a href="http://gnutls.org/security.html">security product</a>,
          it is imperative that you keep your GnuTLS up to date.
        </p>

        <p>
          As with many security products, GnuTLS is maintained by a
          responsible and quick-responding team that takes security
          very seriously.  Benefit from their diligence by keeping
          your GnuTLS up to date.
        </p>

        <p>
          We have received reports of issues with older GnuTLS releases.
          Specifically, version 2.12.20 has problems validating
          certificates under certain conditions. Newer releases have
          addressed memory leaks that were able to take down
          Taskserver.
        </p>
      </đd>
    </td>
  </tr>

  <tr>
    <td>
      <dt>
        <span class="glyphicon glyphicon-question-sign"></span>
        ERROR: Could not find common ancestor for ...<br />
        Client sync key not found.
      </dt>
      <dd>
        <p>
          You skipped the important step of running:
          <pre>$ task sync init</pre>

          This performs an initial upload of your pending tasks, and
          sets up a local sync key, which identifies the last sync
          transaction.
        </p>
      </đd>
    </td>
  </tr>
</table>

<a name="debug"></a>
<h4>Debugging</h4>
<p>
  You may wish to try and debug the problem yourself. You will
  probably not. But if you do, here is how.
</p>

<p>
  Both Taskwarrior and Taskserver have a <code>diagnostics</code>
  command, the purpose of which is to show you relevant
  troubleshooting details. Additionally it will indicate problems
  directly, for example, it will tell you if your cert/key files
  are not readable. The output from <code>diagnostics</code> is
  intended to be included in bug reports, and doing so saves you
  a lot of time, because it's the first thing we'll ask for.
</p>

<pre>$ task diagnostics
...
$ taskd diagnostics
...</pre>

<p>
  Read the output of the <code>diagnostics</code> commands carefully,
  there may be several types of problems mentioned, which need to be
  addressed before going further.
</p>

<p>
  The next step would be to run the server in debug mode. First
  shutdown your Taskserver, then launch it interactively:
</p>

<pre>$ taskdctl stop
...
$ taskd server
...</pre>

<p>
  You can hit <code>Ctrl-C</code> to stop this server. For highly
  verbose output, try this:
</p>

<pre>$ taskd server --debug --debug.tls=2
...</pre>

<p>
  Similarly, Taskwarrior has a verbose debug mode, and debug TLS
  mode:
</p>

<pre>$ task rc.debug=1 rc.debug.tls=2 sync
...</pre>

<a name="help"></a>
<h4>Getting Help</h4>
<p>
  As a last resort, ask for help. But please make sure you have
  carefully reviewed your setup, and gone through the checks above
  before asking. No one wants to lead you through the steps above
  to discover that you didn't.
</p>

<p>
  We'll ask you to provide the <code>diagnostics</code> output
  for both Taskwarrior and Taskserver, then we're going to go
  through the steps above, because this is our checklist also.
</p>

<p>
  Impress us and mention the codeword 'Weasel', to prove that you
  at least read this page all the way to the end.
</p>

<p>
  There are several ways of getting help:

  <ul>
    <li>
      Submit your details to our
      <a href="https://answers.tasktools.org">Q &amp; A site</a>,
      then wait patiently for the community to respond.
    </li>
    <li>
      Email us at
      <a href="mailto:support@taskwarrior.org">support@taskwarrior.org</a>,
      then wait patiently for a volunteer to respond.
    </li>
    <li>
      Join us IRC in the #taskwarrior channel on Freenode.net, and
      get a quick response from the community, where, as you have
      anticipated, we will walk you through the checklist above.
    </li>
  </ul>
</p>

%%%
\fi

\end{document}

% \begin{frame}[fragile]\frametitle{Title}
% \begin{lstlisting}
% \end{lstlisting}
% \end{frame}
%
% \begin{frame}\frametitle{title}
% \begin{center}
% \includegraphics[width=10cm,height=7.5cm]{name.png}
% \end{center}
% \end{frame}
%
% \begin{frame}\frametitle{title}
% \begin{itemize}
% \item \textbf{task {\tt<}filter{\tt>} modify}
% \end{itemize}
% \end{frame}
